%% La date des cendres la plus tôt est le 4 février
%% La date des Rameaux la plus tard est le 18 avril
%% => sanctoral du 5 février au 17 avril

%%%%%%%%%%%%%%%%%%%%%%%%%%%%%%%%%%%%%%%%%%%%%%%%%%%%%%

\festum{3}{0205}{sainte Agathe, vierge et martyre}{5 février}
\saintindex{Agathæ}{Agathe}

\rubrica{Laudes et Vêpres à l'antiphonaire. Petites heures au psautier, avec les capitules, répons brefs et versets du Commun des saintes femmes (pour une martyre), et non du Commun des vierges. Oraison \normaltext{Deus, qui inter cétera} du commun, avec \normaltext{beátæ \rubr{N.} Vírginis et Mártyris tuæ}.}

%%%%%%%%%%%%%%%%%%%%%%%%%%%%%%%%%%%%%%%%%%%%%%%%%%%%%%

\festum{3}{0206}{saint Tite, évêque et confesseur}{6} % mtv 
\saintindex{Titum}{Tite}
\saintindex{Dorótheæ}{Dorothée}

\oratio{Deus, qui beátum Titum Confessórem tuum atque Pontíficem apostólicis virtútibus decorásti: ejus méritis et intercessióne concéde; ut juste et pie vivéntes in hoc sǽculo, ad cæléstem pátriam perveníre mereámur. \perdominumla}{Dieu, qui as orné le bienheureux Tite, ton confesseur et pontife, des vertus apostoliques accorde-nous, par ses mérites et son intercession, que, vivant justement et pieusement en ce monde, nous méritions de parvenir à la céleste patrie. \perdominumfr}

\rubrica{À Laudes, comm. de sainte Dorothée, vierge et martyre, \aarub \scorename{MUVXLAB},% Simile est
\vvrub \normaltext{Diffúsa est}, oraison \normaltext{Indulgéntiam nobis} du commun, avec le nom \normaltext{Doróthea}.}

%%%%%%%%%%%%%%%%%%%%%%%%%%%%%%%%%%%%%%%%%%%%%%%%%%%%%%

\festum{3}{0207}{saint Romuald, abbé}{7} % mtv
\saintindex{Romuáldi}{Romuald}

\rubrica{Oraison \normaltext{Intercéssio nos} du commun, avec le nom \normaltext{Romuáldi}.}

%%%%%%%%%%%%%%%%%%%%%%%%%%%%%%%%%%%%%%%%%%%%%%%%%%%%%%

\festum{3}{0208}{saint Jean de Matha, confesseur}{8 février} % mtv
\saintindex{Joannis de Matha}{Jean de Matha}

\oratio{Deus, qui per sanctum Joánnem órdinem sanctíssimæ Trinitátis ad rediméndum de potestáte Saracenórum captívos cǽlitus institúere dignátus es:\pscross{} præsta, quǽsumus; ut, ejus suffragántibus méritis,\psstar{} a captivitáte córporis et ánimæ, te adjuvánte, liberémur. \perdominumla}{Dieu qui, par saint Jean, as daigné instituer miraculeusement l'Ordre de la Très Sainte Trinité pour racheter les captifs du pouvoir des Sarrasins; nous t'en prions, fais que, par les suffrages de ses mérites, nous soyons délivrés, avec ton secours, de la captivité du corps et de l’âme. \perdominumfr}

%%%%%%%%%%%%%%%%%%%%%%%%%%%%%%%%%%%%%%%%%%%%%%%%%%%%%%

\festum{3}{0209}{saint Cyrille, évêque d'Alexandrie, confesseur et docteur de l'Église}{9}
\saintindex{Cyrílli}{Cyrille}
\saintindex{Apollóniæ}{Apolline}

\oratio{Deus, qui beátum Cyríllum Confessórem tuum atque Pontíficem divínæ maternitátis beatíssimæ Vírginis Maríæ assertórem invíctum effecísti: concéde, ipso intercedénte; ut, qui vere eam Genetrícem Dei crédimus, matérna eíusdem protectióne salvémur. \pereumdemla}{Dieu, qui as fait du bienheureux Cyrille, ton confesseur et pontife, l'invincible défenseur de la divine maternité de la bienheureuse Vierge Marie, accorde-nous par son intercession, à nous qui la croyons vraiment Mère de Dieu, d’être sauvés par sa maternelle protection. \pereumdemfr}

\rubrica{À Laudes, comm. de sainte Apolline, vierge et martyre, \aarub \scorename{MUVXLAB},% Simile est
\vvrub \normaltext{Diffúsa est}, oraison \normaltext{Deus, qui inter cétera} du commun, avec le nom \normaltext{Apollóniæ}.}

\rubrica{À Vêpres, \aarub \scorename{CODOVAM}% O doctor optime
avec le nom \normaltext{beáte Cyrílle}.}

%%%%%%%%%%%%%%%%%%%%%%%%%%%%%%%%%%%%%%%%%%%%%%%%%%%%%%

\festum{3}{0210}{sainte Scholastique, vierge}{10}
\saintindex{Scholasticæ}{Scholastique}
\oratio{Deus, qui ánimam beátæ Vírginis tuæ Scholásticæ ad ostendéndam innocéntiæ viam in colúmbæ spécie cælum penetráre fecísti: da nobis ejus méritis et précibus ita innocénter vívere; ut ad ætérna mereámur gáudia perveníre. \perdominumla}{Dieu qui, pour manifester la voie d’innocence suivie par ta bienheureuse vierge Scholastique, as fait monter au ciel son âme sous la forme d’une colombe, donne-nous, par ses mérites et ses prières, de vivre avec une telle innocence que nous méritions de parvenir aux joies éternelles. \perdominumfr}

%%%%%%%%%%%%%%%%%%%%%%%%%%%%%%%%%%%%%%%%%%%%%%%%%%%%%%

\festum{3}{0211}{Notre-Dame de Lourdes}{11 février}
\saintindex{Mariæ!Apparitio}{Marie!Apparition}

\rubrica{Laudes et Vêpres à l'antiphonaire. Petites heures au psautier, hymne, capitule etc. ci-dessous.}

\oratio{Deus, qui per immaculátam Vírginis Conceptiónem dignum Fílio tuo habitáculum præparásti: súpplices a te quǽsumus; ut, ejúsdem Vírginis Apparitiónem celebrántes, salútem mentis et córporis consequámur. \pereumdemla}{Dieu, qui par l’immaculée Conception de la Vierge, as préparé à ton Fils une habitation digne de lui; nous t'en supplions, fais que, célébrant l’apparition de cette même Vierge, nous obtenions la santé de l’âme et du corps. \pereumdemfr}

\titulumrubrica{À Prime}

\rubrica{Hymne p. \pageref{M-CBMVPH}. Au répons bref, verset de Noël.}

\titulumrubrica{À Tierce}

\rubrica{Hymne p. \pageref{M-CBMVTH}.}
\capitulum{Ct. 2: 13-14}{Surge, amíca mea, speciósa mea, et veni, colúmba mea, in foramínibus petræ, in cavérna macériæ, osténde mihi fáciem tuam, sonet vox tua in áuribus meis.}{Lève-toi, mon amie, ma toute belle, et viens, ma colombe, dans les fentes du rocher, dans les retraites escarpées, que je voie ton visage, que j’entende ta voix !}
\gscore{0211TR}{}{\rrtrans Qu’elle est celle-ci qui monte du désert, enivrée de délices ? \vvtrans Appuyée sur son bien- aimé.}
\versiculus{Vidérunt eam fíliæ Sion.}{Et beatíssimam prædicavérunt.}{Les filles de Sion l’ont vue.}{Et l’ont proclamée bienheureuse.}

\titulumrubrica{À Sexte}

\rubrica{Hymne p. \pageref{M-CBMVSH}.}
\capitulum{Pr. 8: 18-19}{Mecum sunt divítiæ et glória, opes supérbæ et justítia: mélior est enim fructus meus auro et lápide pretióso.}{Avec moi, la richesse et la gloire, fortune durable et juste prospérité. Mon fruit est meilleur que l’or, que l’or fin.}
\gscore{0211SR}{}{\rrtrans Je suis la mère du bel amour. \vvtrans De la crainte, de la science et de la sainte espérance.}
\versiculus{In me omnis grátia viæ et veritátis.}{In me omnis spes vitæ et virtútis.}{En moi, toute grâce pour montrer le chemin et la vérité.}{En moi, toute espérance de vie et de force.}

\titulumrubrica{À None}

\rubrica{Hymne p. \pageref{M-CBMVNH}.}
\capitulum{Si. 24: 18-19}{In me grátia omnis viæ et veritátis, in me omnis spes vitæ et virtútis. Transíte ad me, omnes qui concupíscitis me, et a generatiónibus meis implémini.}{J’ai reçu toute grâce pour montrer le chemin et la vérité. En moi est toute espérance de vie et de force. Venez à moi, vous qui me désirez, rassasiez-vous de mes fruits.}
\gscore{0211NR}{}{\rrtrans Celui qui me trouvera, trouvera la vie. \vvtrans Et il puisera le salut dans le Seigneur.}
\versiculus{Transíte ad me, omnes qui concupíscitis me.}{Et a generatiónibus meis implémini.}{Venez à moi, vous tous qui me désirez.}{Et rassasiez-vous de mes fruits.}

\titulumrubrica{Aux Complies}
\rubrica{Hymne p. \pageref{M-CBMVCH}.}

%%%%%%%%%%%%%%%%%%%%%%%%%%%%%%%%%%%%%%%%%%%%%%%%%%%%%%

\festum{3}{0212}{les Sept Saints fondateurs des Servites de Marie}{12}
\saintindex{Septem Fundatorum Ordinis Servorum B. M. V.}{Sept Saints fondateurs des Servites de Marie}

\titulumrubrica{Aux Laudes}
\capitulum{1 P. 4: 13}{Caríssimi: Communicántes Christi passiónibus gaudéte,\pscross{} ut et in revelatióne glóriæ ejus\psstar{} gaudeátis exsultántes.}{Bien-aimés, dans la mesure où vous communiez aux souffrances du Christ, réjouissez-vous, afin d’être dans la joie et l’allégresse quand sa gloire se révélera.}
\hymnus{0212LH}
\versiculus{Sit memória illórum in benedictióne.}{Et ossa eórum púllulent de loco suo.}{Que leur mémoire soit en bénédiction.}{Et que leurs os soient féconds dans leurs sépulcres.}
\gscore{0212LAB}{B}{Voici comme il est bon, il est doux pour des frères de vivre ensemble et d'être unis !}
\rubrica{Oraison des Vêpres.}

\rubrica{Aux petites heures, tout au commun d'un confesseur non pontife, p. \pageref{M-COPOTR} et suivantes, sauf le capitule de Tierce, \normaltext{Caríssimi}, comme ci-dessous aux Vêpres.}

\titulumrubrica{Aux Vêpres}
\capitulum{1 P. 4: 13}{Caríssimi: Communicántes Christi passiónibus gaudéte,\pscross{} ut et in revelatióne glóriæ ejus\psstar{} gaudeátis exsultántes.}{Bien-aimés, dans la mesure où vous communiez aux souffrances du Christ, réjouissez-vous, afin d’être dans la joie et l’allégresse quand sa gloire se révélera.}
\hymnus{0212VH}
\versiculus{Hi viri misericórdiæ sunt, quorum pietátes non defuérunt.}{Semen eórum et glória eórum non derelinquétur.}{Ceux-ci sont des hommes de miséricorde dont les œuvres de piété n’ont pas fait défaut.}{Ni leur race, ni leur gloire ne seront abandonnées.}
\gscore{0212VAM}{M}{Leur nom demeure éternellement, se perpétuant dans leurs fils qui sont la gloire de ces hommes saints.}
\oratio{Dómine Jesu Christe, qui ad recoléndam memóriam dolórum sanctíssimæ Genetrícis tuæ, per septem beátos Patres nova Servórum ejus família Ecclésiam tuam fecundásti:\pscross{} concéde propítius; ita nos eórum consociári flétibus,\psstar{} ut perfruámur et gáudiis. \quivivisla}{Seigneur Jésus-Christ, pour honorer le souvenir des souffrances de ta très sainte Mère, tu as enrichi ton Église de la nouvelle famille de ses Serviteurs, par sept bienheureux Pères: accorde-nous la faveur de nous associer à leurs larmes pour que nous jouissions de leurs joies. \quivivisfr}

%%%%%%%%%%%%%%%%%%%%%%%%%%%%%%%%%%%%%%%%%%%%%%%%%%%%%%

\festum{4}{0214}{saint Valentin, prêtre et martyr}{14}
\saintindex{Valentini}{Valentin}
\oratio{Præsta, quǽsumus, omnípotens Deus:\pscross{} ut, qui beáti Valentíni, Mártyris tui, natalítia cólimus,\psstar{} a cunctis malis imminéntibus, ejus intercessióne, liberémur. \perdominumla}{Nous t'en prions, Dieu tout-puissant, nous qui honorons la naissance du bienheureux Valentin, ton martyr: fais que nous soyons délivrés, par son intercession, de tous les maux qui nous menacent. \perdominumfr}

%%%%%%%%%%%%%%%%%%%%%%%%%%%%%%%%%%%%%%%%%%%%%%%%%%%%%%

\festum{4}{0215}{saints Faustin et Jovite, martyrs}{15}
\saintindex{Faustíni et Jovítæ}{Faustin et Jovite}

\rubrica{Oraison \normaltext{Deus, qui nos ánnua} du commun, avec les noms \normaltext{Faustíni et Jovítæ}.}

%%%%%%%%%%%%%%%%%%%%%%%%%%%%%%%%%%%%%%%%%%%%%%%%%%%%%%

\festum{4}{0218}{saint Siméon, évêque et martyr}{18}
\saintindex{Simeónis}{Siméon}

\rubrica{Oraison \normaltext{Infirmitátem nostram} du commun, avec le nom \normaltext{Simeónis}.}

%%%%%%%%%%%%%%%%%%%%%%%%%%%%%%%%%%%%%%%%%%%%%%%%%%%%%%

\festum{2}{0222}{Chaire de saint Pierre, apôtre}{22 février}
\saintindex{Petri!Cathedra}{Pierre!Chaire}

\rubrica{Aux petites heures, antiennes et psaumes du jour, le reste du commun d'un confesseur pontife.}
\oratio{Deus, qui beáto Petro Apóstolo tuo, collátis clávibus regni cæléstis, ligándi atque solvéndi pontifícium tradidísti:\pscross{} concéde; ut, intercessiónis ejus auxílio,\psstar{} a peccatórum nostrórum néxibus liberémur: \quivivisla}{Dieu, en confiant au bienheureux Pierre, ton Apôtre, les clefs du royaume céleste, tu lui as donné le pouvoir de lier et de délier: accorde-nous d’être, par le secours de son intercession, libérés des liens de nos péchés. \quivivisfr}

%%%%%%%%%%%%%%%%%%%%%%%%%%%%%%%%%%%%%%%%%%%%%%%%%%%%%%

\festum{3}{0223}{saint Pierre Damien, évêque, confesseur et docteur de l'Église}{23}
\saintindex{Petri Damiani}{Pierre Damien}

\oratio{Concéde nos, quǽsumus, omnípotens Deus: beáti Petri Confessóris tui atque Pontíficis mónita et exémpla sectári;\pscross{} ut per terréstrium rerum contémptum\psstar{}	 ætérna gáudia consequámur. \perdominumla}{Accorde-nous, nous te le demandons, Dieu tout-puissant, de suivre les leçons et les exemples du bienheureux Pierre, ton confesseur et pontife, afin que, par le mépris des choses terrestres, nous obtenions les joies éternelles. \perdominumfr}

\rubrica{À Vêpres, \aarub \scorename{CODOVAM}% O doctor optime
avec le nom \normaltext{beáte Petre}.}

%%%%%%%%%%%%%%%%%%%%%%%%%%%%%%%%%%%%%%%%%%%%%%%%%%%%%%

\festum{2}{0224}{saint Matthias, apôtre}{24 février (le 25 les années bissextiles)}
\saintindex{Matthíæ}{Matthias}

\oratio{Deus, qui beátum Matthíam Apostolórum tuórum collégio sociásti:\pscross{} tríbue, quǽsumus; ut, ejus interventióne,\psstar{} tuæ circa nos pietátis semper víscera sentiámus. \perdominumla}{Dieu, qui as associé le bienheureux Matthias au collège de tes apôtres, nous t'en prions: accorde-nous que, par son intercession, nous ressentions toujours les effets de ta miséricorde. \perdominumfr}

%%%%%%%%%%%%%%%%%%%%%%%%%%%%%%%%%%%%%%%%%%%%%%%%%%%%%%

\festum{3}{0227}{saint Gabriel de la Vierge des douleurs, confesseur}{27 (28)}
\saintindex{Gabrielis a Virgine Perdolente}{Gabriel de la Vierge des douleurs}

\oratio{Deus, qui beátum Gabriélem dulcíssimæ Matris tuæ dolóres assídue recólere docuísti, ac per illam sanctitátis et miraculórum glória sublimásti:\pscross{} da nobis, ejus intercessióne et exémplo; ita Genetrícis tuæ consociári flétibus,\psstar{} ut matérna ejúsdem protectióne salvémur: \quivivisla}{Dieu, qui as appris au bienheureux Gabriel à vénérer assidûment les douleurs de ta très douce Mère, et qui, par elle l’as élevé à la gloire de la sainteté et des miracles, accorde-nous, par son intercession et son exemple, de si bien nous associer aux larmes de ta Mère, que nous soyons sauvés par sa protection. \quivivisfr}

%%%%%%%%%%%%%%%%%%%%%%%%%%%%%%%%%%%%%%%%%%%%%%%%%%%%%%

\festum{0304}{S. Casimiri Confessoris}{4}
\saintindex{Casimiri}{Casimir}
\saintindex{Lúcii}{Lucius}

\oratio{Deus, qui inter regáles delícias et mundi illécebras sanctum Casimírum virtúte constántiæ roborásti:\pscross{} quǽsumus; ut ejus intercessióne fidéles tui terréna despíciant,\psstar{} et ad cæléstia semper aspírent. \perdominumla}{Ô Dieu, qui, parmi les délices royales et les séductions du monde, avez fortifié saint Casimir, par la vertu de constance, faites, nous vous demandons, que, par son intercession, vos fidèles méprisent les biens de la terre et aspirent toujours à ceux du ciel. \perdominumfr}

\rubrica{À Laudes, comm. de saint Lucius, pape et martyr, \aarub \scorename{UMEXLAB},% Qui odit
\vvrub \normaltext{Justus ut palma}, oraison \normaltext{Gregem tuum} du commun, avec le nom \normaltext{Lúcium}.}

%%%%%%%%%%%%%%%%%%%%%%%%%%%%%%%%%%%%%%%%%%%%%%%%%%%%%%

\festum{3}{0306}{saintes Perpétue et Félicité, martyres}{6}
\saintindex{Perpétuæ et Felicitátis}{Perpétue et Félicité}

\rubrica{Tout du commun des saintes femmes, pour plusieurs martyres. Oraison \normaltext{Da nobis} avec les noms \normaltext{Perpétuæ et Felicitátis}.}

%%%%%%%%%%%%%%%%%%%%%%%%%%%%%%%%%%%%%%%%%%%%%%%%%%%%%%

\festum{3}{0307}{saint Thomas d'Aquin, confesseur et docteur de l'Église}{7}
\saintindex{Thomæ de Aquino}{Thomas d'Aquin}

\oratio{Deus, qui Ecclésiam tuam beáti Thomæ Confessóris tui mira eruditióne claríficas, et sancta operatióne fœcúndas: da nobis, quǽsumus; et quæ dócuit, intelléctu conspícere, et quæ egit, imitatióne complére. \perdominumla}{Dieu, qui éclaires ton Église par l'admirable science du bienheureux Thomas, ton confesseur, et qui la fécondes par sa sainte opération; nous t'en prions, fais-nous contempler par l'intelligence ce qu’il a enseigné, et accomplir, en l’imitant, ce qu’il a fait. \perdominumfr}

\rubrica{À Vêpres, \aarub \scorename{CODOVAM}% O doctor optime
avec le nom \normaltext{beáte Thoma}.}

%%%%%%%%%%%%%%%%%%%%%%%%%%%%%%%%%%%%%%%%%%%%%%%%%%%%%%

\festum{3}{0308}{saint Jean de Dieu, confesseur}{8}
\saintindex{Joannis a Deo}{Jean de Dieu}

\oratio{Deus, qui beátum Joánnem, tuo amóre succénsum, inter flammas innóxium incédere fecísti, et per eum Ecclésiam tuam nova prole fecundásti:\pscross{} præsta, ipsíus suffragántibus méritis; ut igne caritátis tuæ vítia nostra curéntur,\pscross{} et remédia nobis ætérna provéniant. \perdominumla}{Dieu, qui as fait marcher indemne au milieu des flammes le bienheureux Jean embrasé de ton amour et qui, par lui, as enrichi ton Église d’une nouvelle famille, fais que, par l’intercession de ses mérites, le feu de ta charité guérisse nos misères, et nous obtienne les remèdes éternels. \perdominumfr}

%%%%%%%%%%%%%%%%%%%%%%%%%%%%%%%%%%%%%%%%%%%%%%%%%%%%%%

\festum{3}{0309}{sainte Françoise Romaine, veuve}{9 mars}
\saintindex{Franciscæ Romanæ}{Françoise Romaine}

\oratio{Deus, qui beátam Francíscam fámulam tuam, inter cétera grátiæ tuæ dona, familiári Angeli consuetúdine decorásti:\pscross{} concéde, quǽsumus; ut, intercessiónis ejus auxílio,\psstar{} Angelórum consórtium cónsequi mereámur. \perdominumla}{Dieu, qui, entre autres dons de ta grâce, as honoré la bienheureuse Françoise, ta servante, de l'habituelle familiarité d’un ange, accorde-nous, nous t'en prions, de mériter par le secours de son intercession, d’être introduits dans la société des anges. \perdominumfr}

%%%%%%%%%% LE 0309 EST LA DATE LA PLUS TARDIVE POSSIBLE POUR LE MARDI AVANT LES CENDRES.

\festum{0310}{Ss. Quadraginta Martyrum}{10 mars}
\saintindex{}
\saintindex{}
\oratio{Præsta, quǽsumus, omnípotens Deus: ut, qui gloriósos Mártyres fortes in sua confessióne cognóvimus, pios apud te in nostra intercessióne sentiámus. \perdominumla}{Faites, nous vous le demandons, Dieu tout-puissant, qu’ayant reconnu la force de ces glorieux Martyrs dans la confession de leur foi, nous sentions l’effet de leur piété dans leur intercession pour nous. \perdominumfr}
\festum{0312}{S. Gregorii Papæ Confessoris et Ecclesiæ Doctoris}{12 mars}
\saintindex{Gregóri}
\saintindex{Gregori}
\oratio{Deus, qui ánimæ fámuli tui Gregórii ætérnæ beatitúdinis prǽmia contulísti: concéde propítius; ut, qui peccatórum nostrórum póndere prémimur, ejus apud te précibus sublevémur. \perdominumla}{Ô Dieu, qui avez accordé à l’âme de votre serviteur Grégoire la récompense de la béatitude éternelle, faites-nous cette miséricorde, qu’accablés sous le poids de nos péchés, nous soyons soulagés par ses prières auprès de vous. \perdominumfr}
\festum{0317}{S. Patricii Episcopi et Confessoris}{17 mars}
\saintindex{Patrícium}
\saintindex{Patrick}
\oratio{Deus, qui ad prædicándam géntibus glóriam tuam beátum Patrícium Confessórem atque Pontíficem míttere dignátus es: ejus méritis et intercessióne concéde; ut, quæ nobis agénda prǽcipis, te miseránte adimplére possímus. \perdominumla}{Ô Dieu, qui avez daigné envoyer le bienheureux Patrick, Confesseur et Pontife, prêcher votre gloire aux nations, accordez-nous, par ses mérites et son intercession, de pouvoir accomplir, grâce à votre miséricorde, ce que vous nous ordonnez de faire. \perdominumfr}
\festum{0318}{S. Cyrilli Episcopi Hierosolymitani Confessoris et Ecclesiæ Doctoris}{18 mars}
\saintindex{Cyrille}
\saintindex{Cyrille}
\oratio{Da nobis, quǽsumus, omnípotens Deus, beato Cyríllo Pontífice intercedénte: te solum verum Deum, et quem misísti Jesum Christum ita cognóscere; ut inter oves, quæ vocem ejus áudiunt, perpétuo connumerári mereámur. \pereumdemla}{Donnez-nous, nous vous le demandons, Dieu tout-puissant, par l’intercession du bienheureux Pontife Cyrille, de vous connaître, vous, le seul vrai Dieu et celui que vous avez envoyé Jésus-Christ, de telle façon que nous méritions d’être comptés éternellement parmi les brebis qui écoutent toujours sa voix. \pereumdemfr}
\festum{0319}{S. Joseph Sponsi B.M.V. Confessoris}{19 mars}
\saintindex{}
\saintindex{}
\oratio{Sanctíssimæ Genetrícis tuæ sponsi, quǽsumus, Dómine, méritis adjuvémur; ut quod possibílitas nostra non óbtinet, ejus nobis intercessióne donétur.
$Qui vivis}{Faites, Seigneur, s’il Vous plaît, que nous soyons secourus par les mérites de l’époux de Votre très sainte Mère afin que ce que nous ne pouvons obtenir de nous-mêmes nous soit donné par son intercession.
$Qui vivis}
\festum{0319}{S. Joseph Sponsi B.M.V. Confessoris}{19 mars}
\saintindex{}
\saintindex{}
\oratio{@Sancti/03-19:Oratio}{}
\festum{0321}{S. Benedicti Abbatis}{21 mars}
\saintindex{Benedícti}
\saintindex{Benoît}
\festum{0323}{S. Turibius of Mogrovejo, bishop}{23 mars}
\saintindex{}
\saintindex{}
\festum{0324}{S. Gabrielis Archangeli}{24 mars}
\saintindex{}
\saintindex{}
\oratio{Deus, qui inter céteros Angelos, ad annuntiándum Incarnatiónis tuæ mystérium, Gabriélem Archángelum elegísti: concéde propítius; ut qui festum ejus celebrámus in terris, ipsíus patrocínium sentiámus in cælis:
$Qui vivis}{Ô Dieu, qui avez choisi l’Archange Gabriel entre tous les Anges, pour annoncer le mystère de Votre Incarnation, accordez-nous, dans Votre bonté, qu’après avoir célébré sa fête sur la terre, nous goûtions dans le ciel les effets de sa protection.
$Qui vivis}
\oratio{Deus, qui per Archángelum Gabriëlem Salvatórem mundi Sacratíssimæ Vírgini concipiéndum nuntiásti: da, ut eúndem et mente pura concipiámus, et férvido imitémur afféctu.
$Qui tecum}{}
\festum{0325}{}{25 mars}
\saintindex{}
\saintindex{}
\oratio{Deus, qui de beátæ Maríæ Vírginis útero Verbum tuum, Angelo nuntiánte, carnem suscípere voluísti: præsta supplícibus tuis; ut, qui vere eam Genetrícem Dei crédimus, ejus apud te intercessiónibus adjuvémur. \pereumdemla}{Ô Dieu, qui avez voulu que votre Verbe prît un corps humain à la parole de l’Ange dans le sein de la bienheureuse Vierge Marie ; accordez à ceux qui vous en supplient que, nous qui la croyons véritablement Mère de Dieu, nous soyons secourus auprès de vous grâce à son intercession. \pereumdemfr}
\festum{0327}{S. Joannis Damasceni Confessoris et Ecclesiæ Doctoris}{27 mars}
\saintindex{Joannes}
\saintindex{Joannes}
\oratio{Omnípotens sempitérne Deus, qui ad cultum sacrárum imáginum asseréndum, beátum Joánnem cælésti doctrína et admirábili spíritus fortitúdine imbuísti: concéde nobis ejus intercessióne et exémplo; ut, quorum cólimus imágines, virtútes imitémur et patrocínia sentiámus. \perdominumla}{Dieu tout-puissant et éternel, qui avez donné au bienheureux Jean une science toute céleste et une admirable force d’âme pour défendre le culte des saintes images, accordez-nous, par son intercession et à son exemple, d’imiter les vertus de ceux dont nous honorons les images, et de ressentir les effets de leur protection. \perdominumfr}
\festum{0328}{S. Joannis a Capistrano Confessoris}{28 mars}
\saintindex{}
\saintindex{}
\oratio{Deus, qui per beátum Joánnem fidéles tuos in virtúte sanctíssimi nóminis Jesu de crucis inimícis triumpháre fecísti: præsta, quǽsumus; ut, spirituálium hóstium, ejus intercessióne, superátis insídiis, corónam justítiæ a te accípere mereámur. \pereumdemla}{Ô Dieu, qui, par le bienheureux Jean, avez fait triompher vos fidèles des ennemis de la croix, grâce à la vertu du très saint nom de Jésus, faites nous vous en supplions, qu’ayant surmonté par son intercession les embûches de nos ennemis spirituels, nous méritions de recevoir de vous la couronne de justice. \pereumdemfr}
\festum{0401}{}{1 avril}
\saintindex{}
\saintindex{}
\oratio{Deus omípotens, cujus sérvitus summa et plena est felícitas: * concéde nobis, intercedénte beáto Hugóne, Abbáte, § in sanctitáte et justítia sine timóre tibi servíre. \perdominumla}{}
\festum{0402}{S. Francisci de Paula Confessoris}{2 avril}
\saintindex{}
\saintindex{}
\oratio{Deus, humílium celsitúdo, qui beátum Francíscum Confessórem Sanctórum tuórum glória sublimásti: tríbue, quǽsumus; ut, ejus méritis et imitatióne, promíssa humílibus prǽmia felíciter consequámur. \perdominumla.}{Ô Dieu, grandeur des humbles, qui avez élevé le bienheureux François, Confesseur, à la gloire de vos Saints, accordez, nous vous en prions, que, par ses mérites et par l’imitation de ses vertus, nous ayons le bonheur d’obtenir les récompenses promises à ceux qui sont humbles. \perdominumfr.}
\festum{0404}{S. Isidori Episcopi Confessoris et Ecclesiæ Doctoris}{4 avril}
\saintindex{Isidórum}
\saintindex{Isidóre}
\oratio{}{Ô Dieu qui avez fait à votre peuple la grâce d’avoir le bienheureux Isidore, pour ministre du salut éternel, faites, nous vous en prions, que nous méritions d’avoir pour intercesseur dans les cieux celui qui nous a donné sur terre la doctrine de vie. \perdominumfr}
\festum{0405}{S. Vincentii Ferrerii Confessoris}{5 avril}
\saintindex{}
\saintindex{}
\oratio{Deus, qui Ecclésiam tuam beáti Vincéntii Confessóris tui méritis et prædicatióne illustráre dignátus es: concéde nobis fámulis tuis; ut et ipsíus instruámur exémplis et ab ómnibus ejus patrocínio liberémur advérsis. \perdominumla}{Dieu, qui avez daigné illustrer Votre Église par Les mérites et la prédication du bienheureux Vincent, Votre confesseur, accordez à nous, Vos familiers, que nous soyons instruits par ses exemples et libérés de toutes adversités par son patronage. \perdominumfr}
\festum{0411}{S. Leonis I Papæ Confessoris et Ecclesiæ Doctoris}{11 avril}
\saintindex{Leónem}
\saintindex{Léon}
\oratio{@Commune/C4-1}{Nous vous supplions, Seigneur, d’exaucer les prières que nous vous adressons en la solennité du bienheureux Léon, votre Confesseur et Pontife, et de nous accorder, grâce aux mérites et à l’intercession de celui qui vous a si dignement servi, le pardon de tous nos péchés. \perdominumfr}
\festum{0413}{S. Hermenegildi Martyris}{13 avril}
\saintindex{}
\saintindex{}
\oratio{Deus, qui beátum Hermenegíldum Mártyrem tuum cælésti regno terrénum postpónere docuísti: da quǽsumus nobis; ejus exémplo cadúca despícere, atque ætérna sectári. \perdominumla}{Ô Dieu, qui avez appris au Bienheureux Herménégilde à faire passer le royaume de ce monde après celui des cieux, accordez-nous, s’il Vous plaît, de mépriser, à son exemple, les biens périssables et de poursuivre les éternels. \perdominumfr}
\festum{0414}{S. Justini Martyris}{14 avril}
\saintindex{}
\saintindex{}
\oratio{Deus, qui per stultítiam crucis eminéntem Jesu Christi sciéntiam beátum Justínum Mártyrem mirabíliter docuísti: ejus nobis intercessióne concéde; ut, errórum circumventióne depúlsa, fídei firmitátem consequámur. \pereumdemla}{Ô Dieu qui, par la folie de la croix, avez merveilleusement enseigné au bienheureux Martyr Justin la suprême science de Jésus-Christ, accordez-nous par Son intercession que, repoussant tout piège d’erreur, nous obtenions la fermeté de la foi. \pereumdemfr}
\festum{0414}{Ss. Tiburtii, Valeriani, et Maximi Martyrum}{14 avril}
\saintindex{}
\saintindex{}
\oratio{Præsta, quǽsumus, omnípotens Deus: ut, qui sanctórum Mártyrum tuórum Tiburtii, Valeriáni et Maximi solemnia colimus; eórum étiam virtútes imitemur. \perdominumla}{Faites, nous Vous en supplions, Seigneur, que nous qui célébrons la fête de Vos saints Martyrs Tiburce, Valérien et Maxime, nous imitions aussi leurs vertus. \perdominumfr}
\festum{0417}{S. Aniceti Papæ et Martyris}{17 avril}
\saintindex{Anicéti}
\saintindex{Anicet}